% !PRINTED 69
\weberpage{69}
Diese Function $\varrho$ gehört dann zu einer gewissen Gruppe $\mathfrak{A}$,
die ein Theiler von $\mathfrak{N}$ ist. Wir nennen auch die Gruppe $\mathfrak{A}$ und
den Körper $\Omega(\varrho)$ \gesperrt{zusammengehörig} und ziehen daraus den
Satz:

\gesperrt{1. Zu jedem Theiler $\mathfrak{A}$ von $\mathfrak{N}$ gehört ein gewisser
in $\Omega_m$ enthaltener Kreistheilungskörper $\Omega(\varrho)$, und
umgekehrt gehört zu jedem Theiler $\Omega(\varrho)$ von $\Omega_m$
ein gewisser Theiler $\mathfrak{A}$ von $\mathfrak{N}$ als Gruppe, in dem
Sinne, dass, wenn $a$ eine Zahl aus $\mathfrak{A}$ ist, alle
Zahlen des Körpers $\Omega(\varrho)$ die Permutationen $(r,r^a)$
gestatten, und dass umgekehrt jede Zahl in $\Omega_m$,
die diese Permutationen gestattet, in $\Omega(\varrho)$ enthalten ist.}

Die Galois'sche Gruppe eines solchen Körpers $\Omega(\varrho)$ erhalten wir nach Bd. I, §. 156, wenn wir $\mathfrak{N}$ in das System der
Nebengruppen zerlegen:
\[
\mathfrak{N} = \mathfrak{A} + \mathfrak{A}_1 + \mathfrak{A}_2 + \cdots,
\]
und die unter den Nebengruppen $\mathfrak{A}, \mathfrak{A}_1, \mathfrak{A}_2 \ldots$ durch Composition
mit den Elementen von $\mathfrak{N}$ hervorgerufenen Permutationen aufsuchen. Diese Gruppe ist aber nach §. 4 isomorph mit der
Gruppe $\mathfrak{N}/\mathfrak{A}$, also auch isomorph mit der zu $\mathfrak{A}$ reciproken Gruppe
(§. 12), die wir mit $\mathfrak{B}$ bezeichnen wollen. Ist $a$ der Grad von $\mathfrak{A}$
und $b$ der von $\mathfrak{B}$, so ist $ab = \nu$, und $\varrho$ genügt einer irreduciblen
Abel'schen Gleichung vom Grade $b$.

\gesperrt{2. Wir bekommen also alle Kreistheilungskörper,
wenn wir zu jedem Modul $m$ die sämmtlichen
Divisoren $\mathfrak{A}$ der Gruppe $\mathfrak{N}$ bilden, zu jeder dieser
Gruppen $\mathfrak{A}$ eine zugehörige Function $\varrho$ suchen
und daraus die Körper $\Omega(\varrho)$ ableiten.}

Es ist aber noch die Frage, ob bei diesem Processe ein und
derselbe Körper $\Omega$ mehrmals auftreten kann, wodurch wir auf
die Untersuchung der gemeinschaftlichen Theiler zweier Kreistheilungskörper geführt werden.

Nach der oben gegebenen Definition ist wohl zu unterscheiden zwischen der Gruppe, zu der ein Körper $\Omega$ \gesperrt{gehört},
und der Galois'schen Gruppe des Körpers; beide Gruppen sind
zu einander reciprok. So gehört der Körper $\Omega_m$ selbst zur
Einheitsgruppe, während seine Galois'sche Gruppe $\mathfrak{N}$ ist,
